\section*{Hoofdstuk \ref{ch:g3:bones-and-muscles} --- Botten en spieren}

\begin{solution}{exo:g3:bones-and-muscles:1}
De \omterm{def:g3:bones-and-muscles:skeleton}{schedel} beschermt de hersenen; de \omterm{def:g3:bones-and-muscles:skeleton}{wervelkolom} houdt de \omterm{def:g1:my-body:parts}{romp} rechtop
(een buigzame zuil); de \omterm{def:g3:bones-and-muscles:skeleton}{ribben} vormen een beschermende kooi om hart en
longen; de pijpbeenderen dragen \omterm{def:g1:my-body:parts}{armen} en \omterm{def:g1:my-body:parts}{benen}.
\end{solution}

\begin{solution}{exo:g3:bones-and-muscles:2}
Ongeveer $200$. Het langste is het dijbeen.
\end{solution}

\begin{solution}{exo:g3:bones-and-muscles:3}
Nee --- een gebogen \omterm{def:g3:bones-and-muscles:skeleton}{bot} is een gebroken \omterm{def:g3:bones-and-muscles:skeleton}{bot}. Al het buigen gebeurt in
de \omterm{prop:g3:bones-and-muscles:joints}{gewrichten}, waar twee \omterm{def:g3:bones-and-muscles:skeleton}{botten} elkaar raken.
\end{solution}

\begin{solution}{exo:g3:bones-and-muscles:4}
Het is een zuil van meer dan $20$ botjes, die elk een paar graden op
hun buren kunnen draaien; de vele kleine buigingen tellen op tot één
grote soepele krul.
\end{solution}

\begin{solution}{exo:g3:bones-and-muscles:5}
Hij wordt korter en dikker, en trekt zo aan het \omterm{def:g3:bones-and-muscles:skeleton}{bot} waaraan hij
vastzit. Geen enkele \omterm{def:g3:bones-and-muscles:muscle}{spier} kan duwen.
\end{solution}

\begin{solution}{exo:g3:bones-and-muscles:6}
Omdat een \omterm{def:g3:bones-and-muscles:muscle}{spier} alleen trekt, kan één \omterm{def:g3:bones-and-muscles:muscle}{spier} het \omterm{def:g1:my-body:joint}{gewricht} maar één kant
op bewegen. Buigen: de voorste \omterm{def:g3:bones-and-muscles:muscle}{spier} trekt samen en trekt de onderarm
omhoog terwijl de achterste ontspant. Strekken: de achterste \omterm{def:g3:bones-and-muscles:muscle}{spier}
trekt samen terwijl de voorste ontspant.
\end{solution}

\begin{solution}{exo:g3:bones-and-muscles:7}
Het levende \omterm{def:g3:bones-and-muscles:skeleton}{bot} bouwt de breuk zelf langzaam weer op; het gips houdt de
stukken alleen stil terwijl het \omterm{def:g3:bones-and-muscles:skeleton}{bot} werkt. Het laat zien dat \omterm{def:g3:bones-and-muscles:skeleton}{bot} levend
materiaal is, gevoed door bloed --- niet dood zoals krijt.
\end{solution}

\begin{solution}{exo:g3:bones-and-muscles:8}
Buigen: de voorkant van de \omterm{def:g1:my-body:parts}{arm} zwelt op en wordt hard (die \omterm{def:g3:bones-and-muscles:muscle}{spier} trekt
samen). Strekken: de voorkant wordt zacht terwijl de achterkant van de
\omterm{def:g1:my-body:parts}{arm} aanspant --- het paar wisselt van rol.
\end{solution}

\begin{solution}{exo:g3:bones-and-muscles:9}
Knieën zijn \omterm{prop:g3:bones-and-muscles:joints}{gewrichten} die gebouwd zijn voor grote vouwen en die door
de sterkste \omterm{def:g3:bones-and-muscles:muscle}{spieren} van het lichaam bewogen worden; de \omterm{def:g3:bones-and-muscles:skeleton}{wervelkolom} is
gebouwd voor zachte buigingen verdeeld over veel botjes. Tillen met
gebogen knieën en een rechte rug geeft het zware werk aan de \omterm{def:g1:my-body:parts}{benen} en
spaart de \omterm{def:g3:bones-and-muscles:skeleton}{wervelkolom}.
\end{solution}

\begin{solution}{exo:g3:bones-and-muscles:10}
Elk touwtje kan, net als elke \omterm{def:g3:bones-and-muscles:muscle}{spier}, alleen trekken. Het antwoord van
de poppenmaker is dat van de bioloog: twee touwtjes per beweging ---
een om de \omterm{def:g1:my-body:parts}{arm} omhoog te trekken, een ander om hem weer omlaag te
trekken --- die om de beurt werken, precies zoals de spierparen van het
lichaam.
\end{solution}
